\documentclass[a4paper,12pt]{extarticle}
\usepackage{graphicx}

% Пакеты для русского языка и кодировки
\usepackage[T2A]{fontenc}
\usepackage[utf8]{inputenc}
\usepackage[russian]{babel}

% Пакеты для математики и теорем
\usepackage{amsmath, amssymb, amsthm}

% Пакеты для графики и гиперссылок
\usepackage{hyperref}

% Настройка полей
\usepackage{geometry}
\geometry{top=2cm, bottom=2cm, left=2.5cm, right=1.5cm}

% Для оформления кода и таблиц
\usepackage{listings}
\usepackage{longtable,booktabs,calc}
\usepackage[table]{xcolor}

% Настройка межстрочного интервала
\usepackage{setspace}
\onehalfspacing

% Отключает нумерацию для part, chapter, section и т.д.
\setcounter{secnumdepth}{-3}

\hypersetup{
    colorlinks=true,
    linkcolor=black,
    urlcolor=blue
}


\begin{document}
\thispagestyle{empty}
\setstretch{1.0}
\begin{center}
{\Large\textbf{Формула итоговой оценки за курс}}
\end{center}

Итоговая оценка складывается из оценок за два раздела курса. Сначала
вычисляется оценка за каждый раздел, затем обе оценки объединяются с равными
весами:

\[
\begin{aligned}
    \text{О}_{\text{курс}} &= 0{,}5\,\text{О}_1 + 0{,}5\,\text{О}_2,\\
    \text{О}_1 &= \text{КР}_1,\\[1em]
    \text{О}_2 &= 0{,}25\,\text{ЛБ}_4 + 0{,}25\,\text{ЛБ}_5\\
        &\quad + 0{,}10\,\text{Активность} + 0{,}40\,\text{КР}_2.
\end{aligned}
\]

Подставив оценки разделов в общую формулу, получаем свёрнутую формулу, в которой
сразу указана доля каждого элемента в итоговой оценке за курс:

\[
\boxed{\text{О}_{\text{курс}} = 0{,}5\,\text{КР}_1
    + 0{,}125\,\text{ЛБ}_4 + 0{,}125\,\text{ЛБ}_5
    + 0{,}05\,\text{Активность} + 0{,}2\,\text{КР}_2}
\]

\paragraph{Обозначения и веса}
\begin{itemize}
    \item \textbf{КР1} --- контрольная работа за раздел 1; её вес в оценке за курс --- 50\%.
    \item \textbf{ЛБ4} и \textbf{ЛБ5} --- оценки за лабораторные работы №4 и №5; вес каждой --- 12{,}5\%.
    \item \textbf{Активность} --- оценка за работу на занятиях; её вес --- 5\%.
    \item \textbf{КР2} --- контрольная работа за раздел 2; её вес --- 20\%.
\end{itemize}

Каждая составляющая оценивается по шкале от 0 до 10. Все веса в свёрнутой
формуле в сумме составляют 100\%. Рассчитанная итоговая оценка за курс
округляется до целого по арифметическому правилу: дробная часть от 0{,}5 и выше
округляется вверх.

\medskip
\noindent\textbf{Пример 1}\par
Пусть \textbf{КР1 = 8}, \textbf{ЛБ4 = 7}, \textbf{ЛБ5 = 9},
\textbf{Активность = 10}, \textbf{КР2 = 6}. Сначала вычислим оценку за раздел 2:
\[
\text{О}_2 = 0{,}25\cdot 7 + 0{,}25\cdot 9 + 0{,}10\cdot 10
    + 0{,}40\cdot 6 = 7{,}4.
\]
Затем найдём итоговую оценку за курс:
\[
\text{О}_{\text{курс}} = 0{,}5\cdot 8 + 0{,}5\cdot 7{,}4
    = 7{,}7.
\]
После округления итоговая оценка за курс равна \textbf{8}.

\medskip
\noindent\textbf{Пример 2. КР1 = 7, оценка за раздел 2 = 0}\par
Пусть \textbf{ЛБ4 = 0}, \textbf{ЛБ5 = 0}, \textbf{Активность = 0} и
\textbf{КР2 = 0}. Тогда:
\[
\text{О}_2 = 0{,}25\cdot 0 + 0{,}25\cdot 0 + 0{,}10\cdot 0
    + 0{,}40\cdot 0 = 0,
\qquad
\text{О}_{\text{курс}} = 0{,}5\cdot 7 + 0{,}5\cdot 0 = 3{,}5.
\]
После округления итоговая оценка за курс равна \textbf{4}.

\medskip
\noindent\textbf{Пример 3. КР1 = 2, оценка за раздел 2 = 4}\par
Пусть \textbf{ЛБ4 = 8}, \textbf{ЛБ5 = 4}, \textbf{Активность = 2} и
\textbf{КР2 = 2}. Тогда:
\[
\begin{aligned}
    \text{О}_2 &= 0{,}25\cdot 8 + 0{,}25\cdot 4 + 0{,}10\cdot 2
        + 0{,}40\cdot 2 = 4,\\
    \text{О}_{\text{курс}} &= 0{,}5\cdot 2 + 0{,}5\cdot 4 = 3.
\end{aligned}
\]
Итоговая оценка за курс после округления равна \textbf{3}.
\end{document}
